\documentclass{article}
\usepackage{graphicx} 
\usepackage[margin=2.5cm]{geometry}
\usepackage{natbib}
\usepackage{tikz}
\usepackage{booktabs}

% Shortcuts %
\newcommand{\agents}{A}
\newcommand{\objects}{H}
\newcommand{\match}{\mu}
\newcommand{\linearorders}{\mathcal{L}}


\newtheorem{example}{Example}
\newtheorem{theorem}{Theorem}
\newtheorem{definition}{Definition}
\newtheorem{observation}{Observation}
\newtheorem{proposition}{Proposition}
\newtheorem{conjecture}{Conjecture}
\newtheorem{corollary}{Corollary}
\newtheorem{lemma}{Lemma}

%TikZ commands
\newcommand*\encircled[1]{\tikz[baseline=(char.base)]{
            \node[shape=circle,draw,inner sep=0.25pt,minimum size=1em] (char) {\vphantom{Ag}#1};}}
\newcommand*\framed[1]{\tikz[baseline=(char.base)]{
            \node[draw,inner sep=0.25pt,minimum size=1em] (char) {\vphantom{Ag}#1};}}
            

\title{Report TER}
\author{Julie Dornat}
\date{April 2026}

\begin{document}

\maketitle

\section{Abstract}
\section{Introduction}

Computational social choice~\citep{brandt2016handbook} studies collective decision-making problems involving multiple agents with potentially conflicting preferences. 
Its goal is to design mechanisms that aggregate individual preferences into a collective decision while satisfying desirable properties such as fairness and efficiency. 
Typical fairness criteria include maxmin and envy-freeness, which aim to capture the extent to which agents are satisfied with the outcome.

A central challenge in computational social choice lies in the elicitation and representation of agents' preferences over a set of possible alternatives or allocations of resources. 
Most existing results assume complete preference information, which limits their applicability in realistic scenarios.
In many practical settings, it is unrealistic to assume that agents can fully specify their preferences, as the number of possible allocations grows combinatorially with the number of resources (\textit{Handbook of Computational Social Choice, Chapter 10, p224}). 
This makes preference elicitation both cognitively demanding for agents and computationally challenging. 
As a result, a significant line of research focuses on developing compact and expressive preference models that mitigate these difficulties.

At the same time, collective decision-making must balance multiple, often conflicting objectives. 
As discussed in the \textit{Handbook of Computational Social Choice, Chapter 12}, fairness criteria such as envy-freeness or maximin may conflict with efficiency requirements such as Pareto efficiency or utilitarian optimality, which aim at maximizing the overall welfare and avoiding loss of resources. 
This results in seeking solution able to handle the tradeoff between fairness and efficiency.

This motivates the study of preference representations that capture only partial information, while still allowing one to reason about fairness and efficiency properties of the resulting outcomes.
In this context, partial preferences provide a natural framework to model incomplete knowledge about agents' preferences. 
As emphasized by \textit{Kenig and Kimelfeld} (2019), incomplete preferences constitute a distinct form of uncertainty.

Rather than assuming fully specified rankings, agents may only reveal limited or partial information. 
As discussed by \textit{Jérôme Lang} (2020), this raises a fundamental question: how can one make reliable collective decisions when the available preference information is incomplete? and when can the eliciting process be stopped to deduce reliable solution?
To address this issue, the notions of possible and necessary outcomes have been introduced, allowing one to reason about decisions that hold under some or all completions of the missing information.

Computational social choice has applications in a wide range of domains, including voting, fair division, matching, multi-agent systems, and information retrieval. 
In this work, we focus in particular on fair division and matching problems.

Fair division concerns the allocation of resources among agents in a fair manner, and applies to both divisible and indivisible goods. 
A classical example in the indivisible setting is the house allocation problem, where each agent is assigned a single house. 
Such settings can also be viewed through the lens of matching theory, where agents and objects form two sets to be matched.

Matching theory provides a rich framework for modeling allocation mechanisms, including concepts such as perfect matching and popular matchings. 

\paragraph{Research Question}

This report investigates the following question:

\begin{quote}
How can we design fair allocation or matching mechanisms when agents' preferences are only partially known?
\end{quote}

More specifically, we explore how notions such as possible and necessary criteria can be used to reason about fairness under incomplete preferences, and how these ideas can be connected to matching frameworks.

% However, many voting schemes (and other social
% choice mechanisms) make stringent assumptions about the preferences provided by
% voters. For instance, it is usually assumed that each voter provides a complete pref-
% erence ranking of all alternatives under consideration.

\section{Preference Models}

Partial preferences aim to represent situations in which agents do not fully specify their preferences over all possible alternatives. 
Instead of a complete description, only partial information is available, which induces uncertainty regarding the evaluation of outcomes with respect to properties such as fairness or efficiency.

From an informational perspective, partial preferences can be seen as a compact representation of agents' preferences. 
This is closely related to the notion of representation and compilation complexity discussed in the \textit{Handbook of Computational Social Choice, Chapter 10, section 10.7}, where the goal is to encode preference information in a concise way while preserving the ability to reason about it.

During our literature review, we identified several families of partial preference models, which differ in the type and amount of information elicited from agents.

\paragraph{(1) Preference Representation Models}

\begin{itemize}
    \item \textbf{Ordinal preferences:} agents provide a ranking over alternatives or subsets of alternatives. Extensions such as the responsive set extension (\textit{Fair Division under Ordinal Preferences:
Computing Envy-Free Allocations of Indivisible Goods} Bouveret et al, Brams et al (2005)) allow one to derive preferences over bundles from rankings over individual items.
    \item \textbf{Cardinal preferences:} agents assign numerical utilities to alternatives, capturing the intensity of preferences.
\end{itemize}

\paragraph{(2) Partial Elicitation Models}

\begin{itemize}
    \item \textbf{Truncated preferences:} agents report only part of their ranking.
    \begin{itemize}
        \item \textbf{Top-$k$ preferences:} only the $k$ most preferred alternatives are reported.
        \item \textbf{Bottom-$k$ preferences:} only the $k$ least preferred alternatives are reported, leaving the rest unranked. In this case, the most preferred alternative may remain unknown.
        \item \textbf{Set ranking queries:} agents provide a ranking into a subset of alternatives.
        \item \textbf{Intervals of cardinal preferences:} agents provide ranges of utility values instead of a single utility.
    \end{itemize}
    
    \item \textbf{Pairwise comparisons:} agents indicate which alternative is preferred between two options. These comparisons are central in majority-based methods and Condorcet-based reasoning (e.g., Condorcet winner, Copeland's rule).

    \item \textbf{Positional queries:} agent provides the alternative she ranks at a given position, leading to the construction of scoring vector. Using these scoring vectors in voting context, Kenig and Kimelfeld (2019) analyze the computational complexity of estimating the probability of an election outcome under uncertainty.

    \item \textbf{Approval and veto preferences:}
    \begin{itemize}
        \item \textbf{Approval preferences:} agents classify alternatives into approved and non-approved sets.
        \item \textbf{Veto preferences:} agents reject a subset of alternatives.
    \end{itemize}

    \item \textbf{Acceptability-based preferences:} agents specify which alternatives are acceptable, often modeled using bipartite graphs in matching problems. For instance, in the marriage problem, some agents may find certain allocations unacceptable. Thus, not all objects are allocable to a given agent. Acceptability therefore amounts to defining a list of preferences by excluding those on which the agent issues a right of veto.

    \item \textbf{Query-based elicitation models:}
    \begin{itemize}
        \item \textbf{Next-best queries:} agents iteratively reveal their most preferred alternative among those the remaining ones. We call it top k queries when we stop the retrieval at rank k.
        \item \textbf{Split queries:} agents partition a subset of alternatives into ordered groups. More precisely, the Chapter 10 of the Handbook of Computational Social Choice explains that agent is asked to partition a subset of alternatives into tow or more subsets and order them such that all alternatives in a higher-ranked subset are preferred to all alternatives in a lower-ranked subset. The k-approval is a special case where two blocks of size k and m-k are created.
        \item \textbf{Choice queries:} agents select their most preferred alternative from a subset.
    \end{itemize}
\end{itemize}

\paragraph{(3) Structured Preference Domains}

\begin{itemize}
    \item \textbf{Single-peaked preferences:} agents have a most preferred option, and their preferences decrease as they move away from this peak. It assumes that there exists an axis on which alternatives are positionned independently of the agents. In the paper "Single-peaked consistency and its complexity" Bruno Escoffier, Jérôme Lang, Meltem ¨Ozt¨urk study whether the preference profiles are single peaked and how to determine the axis based on some agents' preferences.
    \item \textbf{Euclidean preferences:} preferences are derived from distances in a geometric space. This multidimensional space is defined by characteristics used to describe and compare alternatives. The paper "Euclidean preferences" of Anna Bogomolnaia and Jean-François Laslier defines Euclidean preferences, such as the geometric representation as a point in a d-dimensional space. The alternative is judged as good as it is close to the ideal point given by the best preference of an agent. The authors also consider convex preferences, which induce intermediate preferences between two options that agent prefers. They study propositions over the sufficient dimension for representing any Euclidean preference profile.
    \item \textbf{Attribute-based preferences:} agents have different preference profile depending on the attribute considered (politics, economics, etc).
\end{itemize}

\paragraph{(4) Probabilistic Models}

\begin{itemize}
    \item \textbf{Probabilistic preference models:} preferences are drawn from a probability distribution (e.g., impartial culture or Polya-Eggenberger urn models), allowing one to reason under uncertainty.
\end{itemize}

The Chapter 10 of the Handbook of Computational Social Choice, presents most of these preferences models.

These models illustrate different ways of reducing the amount of information required from agents while maintaining the ability to reason about collective decisions. 
Not all these models are equally relevant for the study of fairness under incomplete information, which motivates focusing on specific classes in the remainder of this work.

In practice, the choice of a preference model has a direct impact on both the expressiveness of the representation and the computational complexity of the associated problems. 
It therefore depends on the target application and the type of reasoning required. Partial preference models are widely used in applications such as voting and matching problems (e.g., house allocation, hospital-resident or marriage problems). 
Simpler models reduce the elicitation burden and facilitate computation, while richer models allow for more expressiveness at the cost of increased complexity.

A key advantage of partial preferences is that they reduce the amount of information agents need to communicate, which can improve the reliability of the elicited data while maintaining sufficient structure for decision-making.

\section{Fair Division with Partial Preferences}

Fair division is a combinatorial problem that consists in allocating a set of objects among a set of agents while satisfying certain fairness criteria. 
The nature of the agents' preferences plays a central role in the definition of the problem, as it directly impacts both the expressiveness of the model and the computational complexity of the associated decision procedures.

While early works often assume cardinal preferences, this approach is limited in practice, as agents may not be able to provide precise numerical utilities. 
In many situations, agents can only express qualitative comparisons (e.g., preferring one object to another) or partial information (e.g., intervals of utilities).
This has led to the study of ordinal preferences, where agents provide full rankings instead. 
However, eliciting complete rankings can be impractical in settings with a large number of objects, which motivates the study of partial preferences.

In this setting, agents typically report ordinal preferences over individual goods, from which preferences over bundles must be inferred. 
A key contribution in this direction is due to Brams et al.~(2005), who introduce a pairwise dominance relation to compare bundles based on preferences over individual items. 
This idea is closely related to the responsive set extension, which provides a standard way to lift preferences from individual goods to bundles. 

Such extensions form the semantic foundation for defining envy between agents when only preferences over single items are available. 
Therefore, preferences over single items can be interpreted as a compact representation of preferences over bundles.
Moreover, when preferences over individual items are only partially known, this induced preferences over bundles becomes themselves incomplete, which is precisely what motivates the study of possible and necessary variants of envy-freeness.

\paragraph{Fairness under Incomplete Preferences}

When preferences are incomplete, classical fairness notions such as envy-freeness can no longer be directly evaluated. 
Instead, fairness must be defined with respect to all possible completions of the agents' preferences.

Bouveret et al. (2010) and Aziz et al.~(2024) define the notions of possible and necessary envy-freeness as follows:

\begin{definition}[Envy-Freeness]
An allocation is envy-free if no agent prefers the bundle assigned to another agent over her own bundle.
\end{definition}

\begin{definition}[Possible Envy-Freeness (PEF)]
An allocation is possibly envy-free if there exists at least one completion of the incomplete preferences under which the allocation is envy-free.
\end{definition}

\begin{definition}[Necessary Envy-Freeness (NEF)]
An allocation is necessarily envy-free if it is envy-free under all completions of the incomplete preferences.
In other words, an allocation satisfies NEF if it is envy-free for all additive valuations that are consistent with the agents' ordinal preferences over individual goods.
\end{definition}


\paragraph{Computational Aspects}

The introduction of possible and necessary fairness leads to new computational questions. 
In particular, one may ask whether there exists an allocation satisfying a given fairness notion under incomplete information.

Bouveret et al.~(2010) study fair division under ordinal preferences with incomplete information, assuming that agents report only their preferences over individual goods. 
They show that:
\begin{itemize}
    \item deciding whether a possibly envy-free allocation exists is computationally simple;
    \item deciding whether a necessarily envy-free allocation exists is NP-complete in general;
    \item the problem becomes polynomial-time solvable when the number of agents is restricted to two.
\end{itemize}

These results highlight the intrinsic difficulty of reasoning under incomplete preferences, as necessary fairness requires considering all possible completions of the preference profile.
Necessarily envy-freeness is a particularly strong requirement under incomplete preferences.
It ensures that an allocation remains envy-free regardless of how the missing preference information is completed.

Aziz et al.~(2024) refine and strengthen these results by proving that the problem of deciding the existence of a necessarily envy-free allocation (ExistsNEF) remains NP-complete for any fixed number of agents greater than or equal to three. 
They show that the problem is polynomial-time solvable for two agents, but becomes NP-complete when the number of agents is at least three.

\paragraph{Relaxations of Envy-Freeness} Aziz et al.~(2024) also provide several relaxations of envy-freeness.
One prominent example is envy-freeness up to one item (EF1) and its necessarily envy-free counterpart (NEF1):

\begin{definition}[EF1]
An allocation satisfies EF1 if, for every pair of agents, any envy can be eliminated by removing at most one item from the bundle of the envied agent.
\end{definition}

\begin{definition}[NEF1]
An allocation satisfies necessary EF1 if it satisfies EF1 under all completions of the incomplete preferences. 
\end{definition}

While these notions provide meaningful relaxations of envy-freeness, the computational complexity of determining the existence of such allocations has not been characterized in the literature. 
Aziz et al. (2024) show that a EF1 allocation can be obtained via a round-robin procedure.

\paragraph{Conclusion}
These results highlight a fundamental trade-off between criteria satisfaction and computational complexity in fair division under incomplete preferences. 
While weaker notions such as possible envy-freeness are easier to satisfy and verify, stronger notions such as necessary envy-freeness provide strong guarantees but are computationally much more demanding.

This motivates the study of intermediate notions, such as EF1 and NEF1, which are weaken fairness requirements while preserving meaningful guarantees under incomplete preferences.

In the following, we investigate how similar ideas can be adapted to matching problems, in particular through the notion of popularity under incomplete preferences.


\section{Matching and popularity}

This fair division problem is similar to one-sided matching problems, where only one set (agents) has preferences over the other set (objects) and we try to match.
Hall's theorem is a fundamental result in matching theory that defines the condition for the existence of a perfect matching in a bipartite graph.
In the context of house allocation problems, it can be used to determine whether a perfect matching exists.


As in fairness, matching has several criteria of optimality:
\begin{itemize}
    \item \textbf{Maximum weight:} Maximize the total utility of the agents when agents express their preferences in a cardinal manner. It is equivalent to the utilitarian optimality.
    \item  \textbf{Pareto-optimal matching:}  A matching is Pareto-optimal if there is no other matching in which at least one agent is better off, whilst no agent is worse off.
    \item \textbf{Stable matching:} A matching is stable if there is no blocking pair, i.e, coalition of agents who prefer to be matched together rather than with their current match. This definition suggests to be in a two-sided matching setting.
    \item \textbf{Popular matching:} A matching is popular if there is no other matching that is more popular, i.e., that is preferred by more agents.
\end{itemize}

The notion of envy is unrelated to the notion of stability, as the latter does not evaluate the presence of envy on the objects matched to an agent, but rather the presence of a better match for a pair of agents.

This lead to considering multiple trade-offs between the satisfaction of optimality criteria and the number of agents matched. 


\paragraph{Research Questions}

Several questions have been investigated in the matching literature, such as:
\begin{itemize}
    \item How to ensure matching criteria such as Maximum weight or minimum cost matching when we have cardinal preferences?
    \item How to ensure stability when we have ordinal preferences?
\end{itemize}

\paragraph{Algorithms and Complexity}

Rank maximal matching: Matching that maximize the number of agents who
are matched with their most preferred object. More agents are matched to their
first choice, better it is. Can be found in polynomial time


\section{Preliminaries}

A House allocation instance $I$ with partial preferences is given by $<\agents, \objects, P>$ where $\agents=\{a_1, ..., a_n\}$ is a set of $n$ applicants and $\objects = \{h_1, ..., h_n\}$ is a set of $n$ houses. 
There are as many houses as applicants ($|\agents| = |\objects|$) and a house cannot be divided to be allocated to several applicants.

The partial preference profile $P=\{P_1, ..., P_n\}$ is the collection of the lists of incomplete preferences for all applicants. 

Each applicant $a_i$ expresses their incomplete preferences over the set of houses $\objects$ by giving a strict partial order $P_i$ over $\objects$. 
More precisely, $h_l P_i h_k$ denotes the pairwise comparison of the houses $h_l$ and $h_k$ for applicant $a_i$, yielding that applicant $a_i$ prefers the house $h_l$ to $h_k$.
The preferences are transitive, i.e if an applicant $a_i$ prefers a house $h_l$ to a house $h_k$ and $h_k$ is preferred to a house $h_z$, then she prefers $h_l$ to $h_z$.

We assume that, even though we do not observe them because of the elicitation process, each applicant $a_i \in \agents$ has strict ordinal preferences on the houses, represented by a linear order $\succ_i$ over $\objects$.
The complete preference profile $\succ$ is the collection of all applicants' preferences, i.e., $\succ=(\succ_1,\dots,\succ_n)$.

To reconstruct the complete preferences, we will consider possible completions of the partial preference profile $P$.
A linear order $\succ_i$ is a linear extension of $P_i$ if it is compatible with $P_i$, i.e., for all houses $h_l,h_k\in \objects$ if $h_l P_i h_k$ then $h_l \succ_i h_k$. 
A preference profile $\succ$ is a possible completion of a partial preference profile $P$ if $\succ_i$ is a linear extension of $P_i$, for every $a_i\in \agents$.

A matching is a function that assigns houses to applicants. 
Each applicant can be assigned to exactly one house and each house can be assigned to at most one applicant. 
A perfect matching $\match:\agents\to \objects$ is a bijection assigning to every applicant $a_i\in \agents$ a house in $\objects$ denoted by $\match(a_i)$.

A matching $\match$ is \emph{Pareto-optimal} if there is no other matching that is at least as good for all applicants and strictly better for at least one applicant.
We are interested in a refinement of Pareto-optimality called \emph{popularity}: a matching is popular if there is no other matching that is more popular, i.e., that is preferred by more applicants.

\begin{definition}[Popularity]
A matching $\match$ is popular if $|\{i\in \agents: \match(i)\succ_i \match'(i)\}|\geq |\{i\in \agents: \match'(i)\succ_i \match(i)\}|$, for every matching $\match'$.
\end{definition}

A popular matching is trivially Pareto-optimal but is not guaranteed to exist.
The following characterization of popularity due to \citet{abraham2007popular} allows to decide in polynomial time whether a popular matching exists.

\begin{lemma}[\citet{abraham2007popular}]\label{lem:charact-pop}
A matching is popular iff every applicant is assigned either her most preferred item or her most preferred item among those that are not ranked first by any applicant.
\end{lemma}

The definition of popularity can then be adapted to partial preferences by considering possible or necessary outcomes with respect to the possible preference completions.

\begin{definition}[Possible popularity]
A matching is possibly popular if it is popular in at least one completion of the partial preference profile.
\end{definition}

\begin{definition}[Necessary popularity]
A matching is necessarily popular if it is popular in all completions of the partial preference profile. 
\end{definition}

We can then investigate the following computational problems:
\smallbreak
\begin{tabular}{ll}
\toprule
\multicolumn{2}{c}{\textsc{ExistPossibPop}/\textsc{ExistNecessPop}} \\
\midrule
Instance: & House allocation instance with incomplete preferences $(N,P)$ \\
Question: & Does there exist a possibly/necessarily popular matching? \\
\bottomrule
\end{tabular}

\smallbreak
\begin{tabular}{ll}
\toprule
\multicolumn{2}{c}{\textsc{VerifPossibPop}/\textsc{VerifNecessPop}} \\
\midrule
Instance: & House allocation instance with incomplete preferences $(N,P)$, matching $\match$ \\
Question: & Is $\match$ possibly/necessarily popular? \\
\bottomrule
\end{tabular}

\section{Openings}

\bibliographystyle{plainnat}
\bibliography{ref}

\end{document}

